\subsection{Generation Overview}
\label{sec:generation-overview}

A mechanism is the equation executed at every generated step. Its structure,
coefficients, and lag assignments remain fixed while the initial history and
innovation sequence determine a particular trajectory. The same final equation,
initial history, and numerical innovations produce the same trajectory, up to
floating-point effects. A construction method can change that equation before
execution. The original and transformed equations must therefore be retained
separately, with ground-truth explanations referring to the executed mechanism.

Tables~\ref{tab:generation-overview-linear} and~\ref{tab:generation-overview-nonlinear}
show nine sampled examples, selected under restricted sampling configurations
to make these changes readable. The first column gives the original recurrence;
the middle columns give the recurrence produced by each construction; the last
column explains the feature being illustrated. The cases cover individual lags,
an offset, multiple memory scales, repeated factors, and nonlinear interactions.
They are worked examples rather than estimates of how often those structures
occur under the benchmark's full sampling prior.

\paragraph{Reading the equations and the innovations.}
To keep the tables compact, write $z_\ell=x_{t-\ell}$ and
$T_\ell=\tanh(z_\ell)$. A squared lag is displayed as $z_1^2$; its AST contains
two lag leaves joined by multiplication. The index $t$ is the current generated
step, so the first update after history ending at time zero uses $t=1$.
Every recurrence explicitly includes a fresh additive innovation, identified
by its superscript: O denotes the original AST, J the Jacobian method,
F factor saturation, P product-output saturation, and H Heaps. This term is
the generator's random innovation; forecast errors are measured later by
comparing predictions with generated observations.

Here the configured standard deviation is $\sigma_\varepsilon=0.05$.
The O, F, P, and H policies use i.i.d. Gaussian innovations with distribution
$\mathcal N(0,0.05^2)$. The J policy uses i.i.d. uniform innovations on
$[-\sqrt{3}(0.05),\sqrt{3}(0.05)]$ to respect its finite noise bound while
matching that standard deviation. These method labels identify policies;
actual noise sequences may be shared or independently drawn according to the
comparison protocol. Matching standard deviations alone leaves different noise
laws. A trajectory comparison must also specify its initial history and which
numerical innovations are shared. The constructors scale only the deterministic
expression; the innovation is added afterwards at every step.

\paragraph{Construction settings used in the tables.}
The Jacobian calculation uses $r_D=1$, $q_w=1.25$, target gain 0.9, and safety
multiplier 0.95. It scales the complete AST to meet both the sensitivity and
domain-containment checks of Section~\ref{sec:implemented-domain-contraction}.
The two Shao--Wu-inspired modes use target gain 0.8 and safety multiplier 0.95.
Factor saturation replaces lag leaves inside nonlinear products by their
hyperbolic tangents, then checks global sensitivity using the condition of
\citet[Theorem~5.1]{shao2007spectral}. Product-output saturation bounds each
complete nonlinear product and checks the project's second-moment envelope,
as detailed in Section~\ref{sec:implemented-shao-wu}. Their verification records
therefore certify different properties. Both preserve standalone linear terms
before applying their root multiplier. Even a candidate already within a
target bound receives the implementation's 0.95 safety multiplier.

\paragraph{How the Heaps column is matched to a linear example.}
Heaps constructs a linear VAR from free parameter matrices and an innovation
covariance \citep[Sections~2.1--2.2 and Appendix]{heaps2023stationarity}.
For this illustration, each scalar coefficient in the original linear formula
is supplied as the corresponding one-by-one free matrix $\mathbf U_\ell$;
missing lags receive zero. This is an explicit comparison convention, with
the same numerical entries interpreted in different parameter spaces.
The constructor then reconstructs the final AR coefficients jointly. We choose
its stationary mean to preserve the original intercept: after reconstruction,
divide that intercept by one minus the sum of the final AR coefficients.
With one lag, free entry 2 consequently produces coefficient
$2/\sqrt{5}\approx0.894427$. In R4, the free entries at lags 1 and 3 also
produce a coefficient at lag 2. Nonlinear rows indicate the linear scope of
this constructor. Section~\ref{sec:implemented-heaps} explains its standard
stationary-history sampling; the equations below describe the transition rule
separately from that initialization choice.

\paragraph{Reproduction and interpretation.}
The input equations are actual outputs of
\path{src/tsfmxai/generators/random_expression.py}, with the seeds shown in the
tables. Their configurations use one series, no time terminals, and no automatic
lag-coefficient gates; selected parameter magnitudes are fixed at 0.5 or 2 to
keep the arithmetic readable. Operator counts, sampling weights, and lag ranges
for each case are recorded in \path{scripts/check_generation_overview.py}.
Running \texttt{python} \path{scripts/check_generation_overview.py} reconstructs the
nine samples, checks all 27 AST transformations and five matched Heaps
constructions, and verifies the formulas numerically. Coefficients shown with
$\approx$ are rounded to six decimals; calculations retain full precision.
These examples compare construction rules. Explosion rates, initialization
effects, and generation speed require the separate trajectory experiments.

\begin{landscape}
\begingroup
\renewcommand{\baselinestretch}{1}
\footnotesize
\setlength{\tabcolsep}{3pt}
\setlength{\LTcapwidth}{.97\linewidth}
\renewcommand{\arraystretch}{1.2}
\begin{longtable}{@{}
>{\raggedright\arraybackslash}p{.13\linewidth}
>{\raggedright\arraybackslash}p{.14\linewidth}
>{\raggedright\arraybackslash}p{.15\linewidth}
>{\raggedright\arraybackslash}p{.17\linewidth}
>{\raggedright\arraybackslash}p{.17\linewidth}
>{\raggedright\arraybackslash}p{.19\linewidth}@{}}
\caption{Generation Overview: linear and affine sampled recurrences. Each cell
includes the innovation added at step $t$. The Heaps column uses the explicitly
declared free-parameter convention.}\label{tab:generation-overview-linear}\\
\toprule
Sampled equation / original AST & Jacobian: scaled AST & Shao--Wu: factors & Shao--Wu: product output & Heaps: linear construction & Why this example?\\
\midrule
\endfirsthead
\toprule
Sampled equation / original AST & Jacobian: scaled AST & Shao--Wu: factors & Shao--Wu: product output & Heaps: linear construction & Why this example?\\
\midrule
\endhead
\textbf{Operation}\newline
Keep the sampled formula and add the innovation at each step. &
Scale the complete deterministic expression to meet contraction and domain
bounds. In R2, multiply by $0.4275$. &
Apply $\tanh$ to lag factors inside nonlinear products, then scale the whole
expression to limit sensitivity. In R2, multiply by $0.38$. &
Apply $\tanh$ after the complete nonlinear product, then scale the whole
expression to satisfy the second-moment bound. &
Map free parameters to stable AR coefficients. For AR(1),
$a=u/\sqrt{1+u^2}$; $u=2$ gives $a\approx0.894427$. &
Each following row applies these operations to one input equation.
The innovation is added after the deterministic transformation.\\
\midrule
\textbf{R1; seed 0}\newline
$x_t=z_1+\varepsilon_t^{(\mathrm O)}$ &
$x_t=0.855z_1$\newline $+\varepsilon_t^{(\mathrm J)}$ &
$x_t=0.76z_1$\newline $+\varepsilon_t^{(\mathrm F)}$ &
$x_t=0.76z_1$\newline $+\varepsilon_t^{(\mathrm P)}$ &
$x_t\approx0.707107z_1$\newline $+\varepsilon_t^{(\mathrm H)}$ &
Unit feedback: the original is a random walk. All three constructions change
its memory strength.\\
\addlinespace
\textbf{R2; seed 0}\newline
$x_t=2z_1+\varepsilon_t^{(\mathrm O)}$ &
$x_t=0.855z_1$\newline $+\varepsilon_t^{(\mathrm J)}$ &
$x_t=0.76z_1$\newline $+\varepsilon_t^{(\mathrm F)}$ &
$x_t=0.76z_1$\newline $+\varepsilon_t^{(\mathrm P)}$ &
$x_t\approx0.894427z_1$\newline $+\varepsilon_t^{(\mathrm H)}$ &
Amplifying feedback. Scaling maps R1 and R2 to identical J/F/P drifts, while
the Heaps parameter map retains a difference.\\
\addlinespace
\textbf{R3; seed 2}\newline
$x_t=z_1-0.5$\newline $+\varepsilon_t^{(\mathrm O)}$ &
$x_t\approx0.578485$\newline $\times(z_1-0.5)$\newline $+\varepsilon_t^{(\mathrm J)}$ &
$x_t=0.76(z_1-0.5)$\newline $+\varepsilon_t^{(\mathrm F)}$ &
$x_t=0.76(z_1-0.5)$\newline $+\varepsilon_t^{(\mathrm P)}$ &
$x_t\approx0.707107z_1$\newline $-0.5+\varepsilon_t^{(\mathrm H)}$ &
Offset: the AST root multiplier scales the intercept too. J reserves space
for offset and noise; H preserves the intercept by the stated convention.\\
\addlinespace
\textbf{R4; seed 0}\newline
$x_t=z_1+z_3$\newline $+\varepsilon_t^{(\mathrm O)}$ &
$x_t\approx0.333659$\newline $\times(z_1+z_3)$\newline $+\varepsilon_t^{(\mathrm J)}$ &
$x_t=0.38(z_1+z_3)$\newline $+\varepsilon_t^{(\mathrm F)}$ &
$x_t=0.38(z_1+z_3)$\newline $+\varepsilon_t^{(\mathrm P)}$ &
$x_t\approx0.707107z_1$\newline $-0.5z_2+0.707107z_3$\newline $+\varepsilon_t^{(\mathrm H)}$ &
Separated lags: J weights older memory more strongly. Joint H reconstruction
introduces a lag-2 coefficient.\\
\addlinespace
\textbf{R5; seed 0}\newline
$x_t=z_{12}+\varepsilon_t^{(\mathrm O)}$ &
$x_t\approx0.073444z_{12}$\newline $+\varepsilon_t^{(\mathrm J)}$ &
$x_t=0.76z_{12}$\newline $+\varepsilon_t^{(\mathrm F)}$ &
$x_t=0.76z_{12}$\newline $+\varepsilon_t^{(\mathrm P)}$ &
$x_t\approx0.707107z_{12}$\newline $+\varepsilon_t^{(\mathrm H)}$ &
Long delay: the J weight $1.25^{11}$ requires substantially more attenuation
than in R1. Formula size alone misses this difficulty.\\
\bottomrule
\end{longtable}
\endgroup
\end{landscape}

\begin{landscape}
\begingroup
\renewcommand{\baselinestretch}{1}
\footnotesize
\setlength{\tabcolsep}{3pt}
\setlength{\LTcapwidth}{.97\linewidth}
\renewcommand{\arraystretch}{1.2}
\begin{longtable}{@{}
>{\raggedright\arraybackslash}p{.13\linewidth}
>{\raggedright\arraybackslash}p{.14\linewidth}
>{\raggedright\arraybackslash}p{.15\linewidth}
>{\raggedright\arraybackslash}p{.17\linewidth}
>{\raggedright\arraybackslash}p{.17\linewidth}
>{\raggedright\arraybackslash}p{.19\linewidth}@{}}
\caption{Generation Overview: nonlinear sampled recurrences.
Here $z_\ell=x_{t-\ell}$ and $T_\ell=\tanh(z_\ell)$.
Factor mode verifies global sensitivity; product-output mode verifies a
second-moment bound.}\label{tab:generation-overview-nonlinear}\\
\toprule
Sampled equation / original AST & Jacobian: scaled AST & Shao--Wu: factors & Shao--Wu: product output & Heaps: linear construction & Why this example?\\
\midrule
\endfirsthead
\toprule
Sampled equation / original AST & Jacobian: scaled AST & Shao--Wu: factors & Shao--Wu: product output & Heaps: linear construction & Why this example?\\
\midrule
\endhead
\textbf{Operation}\newline
Keep the sampled formula and add the innovation at each step. &
Scale the complete deterministic expression to meet contraction and domain
bounds. In R2, multiply by $0.4275$. &
Apply $\tanh$ to lag factors inside nonlinear products, then scale the whole
expression to limit sensitivity. In R2, multiply by $0.38$. &
Apply $\tanh$ after the complete nonlinear product, then scale the whole
expression to satisfy the second-moment bound. &
Map free parameters to stable AR coefficients. For AR(1),
$a=u/\sqrt{1+u^2}$; $u=2$ gives $a\approx0.894427$. &
Each following row applies these operations to one input equation.
The innovation is added after the deterministic transformation.\\
\midrule
\textbf{R6; seed 0}\newline
$x_t=z_1z_2+\varepsilon_t^{(\mathrm O)}$ &
$x_t=0.38z_1z_2$\newline $+\varepsilon_t^{(\mathrm J)}$ &
$x_t=0.38T_1T_2$\newline $+\varepsilon_t^{(\mathrm F)}$ &
$x_t=0.95\tanh(z_1z_2)$\newline $+\varepsilon_t^{(\mathrm P)}$ &
Outside the linear VAR vocabulary. &
Two distinct lag factors: contrast a scaled polynomial, bounded factors,
and a bounded complete product.\\
\addlinespace
\textbf{R7; seed 0}\newline
$x_t=z_1^2+\varepsilon_t^{(\mathrm O)}$ &
$x_t=0.4275z_1^2$\newline $+\varepsilon_t^{(\mathrm J)}$ &
$x_t=0.38T_1^2$\newline $+\varepsilon_t^{(\mathrm F)}$ &
$x_t=0.95\tanh(z_1^2)$\newline $+\varepsilon_t^{(\mathrm P)}$ &
Outside the linear VAR vocabulary. &
Repeated coordinate: two leaves refer to one lag. The conservative
derivative bound counts both occurrences.\\
\addlinespace
\textbf{R8; seed 5}\newline
$x_t=z_1^2z_2z_3$\newline $+\varepsilon_t^{(\mathrm O)}$ &
$x_t\approx0.177662$\newline $\times z_1^2z_2z_3$\newline $+\varepsilon_t^{(\mathrm J)}$ &
$x_t=0.19T_1^2T_2T_3$\newline $+\varepsilon_t^{(\mathrm F)}$ &
$x_t=0.95$\newline $\times\tanh(z_1^2z_2z_3)$\newline $+\varepsilon_t^{(\mathrm P)}$ &
Outside the linear VAR vocabulary. &
Four factors: increasing product degree raises the sensitivity budget.
P saturates the entire connected product once.\\
\addlinespace
\textbf{R9; seed 1438}\newline
$x_t=0.5z_1+z_1z_2$\newline $+\varepsilon_t^{(\mathrm O)}$ &
$x_t\approx0.310909$\newline $\times(0.5z_1+z_1z_2)$\newline $+\varepsilon_t^{(\mathrm J)}$ &
$x_t=0.304$\newline $\times(0.5z_1+T_1T_2)$\newline $+\varepsilon_t^{(\mathrm F)}$ &
$x_t=0.95$\newline $\times[0.5z_1+\tanh(z_1z_2)]$\newline $+\varepsilon_t^{(\mathrm P)}$ &
Outside the linear VAR vocabulary. &
Mixed linear and nonlinear feedback: selective saturation preserves the
linear lag; the root multiplier still scales both terms.\\
\bottomrule
\end{longtable}
\endgroup
\end{landscape}
